\begin{verbatim}
    class Depreciation:
        def __init__(self, amt, year, rate):
            self.amt = amt
            self.year = year
            self.rate = rate
    
        def calculate(self):
            value = self.amt
            for _ in range(self.year):
                value -= value * self.rate / 100
            return value
    
    amt = float(input("Enter purchase value: "))
    year = int(input("Enter years of service: "))
    rate = float(input("Enter depreciation rate (%): "))
    
    asset = Depreciation(amt, year, rate)
    
    print("Depreciated value =", asset.calculate())
\end{verbatim}